\cleardoublepage
\newrefsection
\chapter{文献综述}

\section{MEV与区块链机器人研究综述}

\subsection{MEV概念的提出与演化}

\par 最大可提取价值（Maximal Extractable Value, MEV）的概念最早由Daian等人在2019年的论文《Flash Boys 2.0》中系统阐述，其前身为“矿工可提取价值”（Miner Extractable Value）。MEV指区块生产者（矿工或验证者）通过在其生产的区块内包含、排除或重新排序交易，从交易流中提取的超出标准区块奖励和 gas 费用的价值。MEV的产生源于去中心化金融（DeFi）协议中的套利机会、清算事件以及用户交易被抢跑（front-running）等市场摩擦。
\par 随着以太坊从工作量证明（PoW）向权益证明（PoS）的过渡，以及提议者-构建者分离（Proposer-Builder Separation, PBS）机制的引入，MEV的提取主体从矿工转变为验证者、构建者和搜索者（searchers）组成的复杂生态系统。这一演变改变了MEV市场的结构，也催生了更多自动化参与者的出现。

\subsection{以太坊MEV生态研究}

\par 以太坊作为MEV研究的发源地，其生态体系最为复杂。搜索者（通常为机器人）通过监控内存池（mempool）中的待处理交易，识别套利或清算机会，并提交高 gas 费用的交易以抢占先机。矿工或验证者则通过打包这些交易获得额外收益。
\par 现有研究对以太坊MEV的分类主要包括：
\begin{enumerate}
    \item 套利（Arbitrage）：利用不同DEX间的价格差进行低买高卖。
    \item 三明治攻击（Sandwich Attack）：在用户的大额交易前后插入买卖订单，通过影响价格获利。
    \item 清算（Liquidation）：在借贷协议中，当抵押品价值不足时，机器人抢先清算以获得清算奖励。
\end{enumerate}
\par Qin等人（2021）首次对以太坊MEV进行了量化分析，发现MEV在以太坊早期阶段就已达到数亿美元规模。此后，随着Flashbots等基础设施的出现，MEV的提取逐渐中心化，但也为研究者提供了更透明的数据。\cite {123}

\subsection{BSC上的机器人研究}
\par BSC因其低手续费和高吞吐量，吸引了大量DeFi项目和散户投资者，同时也成为恶意机器人的温床。一个研究团队（Chronis等）连续发表了三篇论文，对BSC上的代币生态和狙击机器人行为进行了系统分析。
\begin{enumerate}
    \item 《Token Spammers, Rug Pulls, and Sniper Bots: An Analysis of the Ecosystem of Tokens in Ethereum and in the Binance Smart Chain (BNB)》（Usenix’23）
    \par 该论文首次对以太坊和BSC上的ERC-20/BEP-20代币生态进行了大规模实证研究。作者构建了包含数百万代币和流动性池的数据集，分析了代币生命周期、创建者行为以及欺诈模式。研究发现：
	\begin{itemize}
  	    \item 约60\%的代币生命周期不足一天，被称为“一日代币”（1-day tokens）。
	    \item 仅1\%的地址创建了超过20\%的代币，这些地址被称为“代币垃圾发送者”（Token Spammers）。
             \item 狙击机器人（Sniper Bot）在新代币上线瞬间抢先购买，频繁参与Rug Pulls交易，既成为诈骗受害者，也助推了骗局的成功。
	\end{itemize}
    \par 该研究的创新之处在于首次系统分析了BSC上的代币生态，但未考虑多地址协同或身份关联的欺诈行为。
    \item 《Ready, Aim, Snipe! Analysis of Sniper Bots and their Impact on the DeFi Ecosystem》（WWW’23）
    \par 该论文聚焦于狙击机器人的行为模式与影响。作者从GitHub收集了28个开源狙击机器人项目，通过代码分析了解其工作流程，并构建以太坊和BSC的流动性池数据集，识别狙击行为。主要结论：
	\begin{itemize}
  	    \item 狙击机器人在BSC上更活跃，策略为“高频小额”；以太坊因高手续费，策略为“低频大额”。
	    \item 以太坊上狙击成功率25.6\%，BSC上仅7\%，大部分失败源于未能及时卖出。
             \item 约17.9\%（以太坊）和37.36\%（BSC）的代币合约至少采用了一种反机器人机制。
	\end{itemize}
    \par 该研究首次提出基于行为模式（操作速度和广度）识别“系列性”狙击机器人地址的检测方法，但样本仅限于开源机器人，忽略了闭源和服务化产品。
    \item 《The Blockchain Warfare: Investigating the Ecosystem of Sniper Bots on Ethereum and BNB Smart Chain》（ACM TOIT 2024）
    \par 作为前述工作的延续，该论文收集了更大规模的数据（至2023年2月），并更细致地分析了交易距离和gas策略。研究发现：
	\begin{itemize}
  	    \item BSC上操作数量是以太坊的5倍，但以太坊成功率更高。
	    \item 狙击机器人在以太坊和BSC上分别投入约1.55亿和1.37亿美元，导致代币价格中位数上涨19.7\%（以太坊）和4.4\%（BSC）。
             \item 以太坊上56.1\%的操作盈利，BSC上为37.5\%。
             \item 分析了近百万智能合约，发现34.4\%（以太坊）和36.2\%（BSC）的合约至少采用一种反机器人机制。
	\end{itemize}
    \par 该研究的创新在于首次系统分析狙击机器人的全生命周期，并量化了反狙击机制的普及率，但识别反制代码的方法依赖预定义正则表达式，可能遗漏新型防御机制。
\end{enumerate}
\par 这三篇论文共同奠定了对EVM链上狙击机器人行为的理解基础，其研究方法（数据集构建、行为特征提取、机器学习检测）对本研究具有重要借鉴意义。

\subsection{Solana MEV研究现状与挑战}

\par 与以太坊和BSC相比，Solana的MEV研究尚处于起步阶段。Solana的独特架构——并行执行、历史证明（PoH）、无公开内存池（mempool）以及交易转发机制——使得MEV的提取方式与EVM链有显著差异。
\begin{itemize}
    \item 无公开内存池：Solana的交易在被确认前不会广播给全网节点，这理论上减少了抢跑（front-running）的可能性。然而，验证者仍可通过提前看到交易流或与交易者合谋来提取价值。
    \item 交易调度权：Solana的领导者（leader）在每轮中拥有交易排序权，可能通过优化排序获利。
    \item Jito Labs：Solana生态中出现了类似Flashbots的MEV基础设施，如Jito Labs，允许搜索者通过小费激励验证者包含其交易，从而创造了一个竞拍市场。
\end{itemize}
\par 现有对Solana MEV的学术研究极少。一些社区报告和博客文章指出，Solana上存在套利机器人和狙击机器人，尤其是在Meme币发行初期。例如，狙击机器人通过监控新代币的创建交易，在流动性池建立瞬间抢先购买，与BSC上的行为类似。
\par 然而，由于缺乏公开内存池和详细的交易元数据，Solana机器人的检测更具挑战性。相关研究者必须依赖链上已确认的交易，通过行为模式而非交易池数据来识别机器人。

\section{区块链机器人检测方法研究}

\par 随着区块链生态中自动化程序的普及，如何从链上海量交易数据中准确识别机器人账户，成为学术界和工业界关注的热点。检测方法经历了从简单的启发式规则到复杂的机器学习模型的演进，并在特征工程、模型选择和验证方式上不断深化。本节系统梳理现有的区块链机器人检测方法，并重点评述与本研究密切相关的经典工作。

\subsection{基于规则的启发式检测方法}

\par 早期研究多依赖领域知识设定阈值规则，通过观察机器人行为与人类行为的差异进行识别。这类方法的优势在于解释性强、计算开销小，但规则设计依赖专家经验，难以覆盖多样化的机器人策略，且阈值选择具有主观性。
\par 常见的启发式规则包括：
\begin{enumerate}
    \item 交易频率规则：机器人通常具有高频交易特征，如平均交易间隔小于10秒、单日交易数超过一定阈值。 
    \item 时间分布规则：机器人不受作息限制，可能在夜间（如UTC 0:00-5:00）或周末保持活跃，而人类用户在这些时段交易量显著下降。
    \item Gas费用规则：在以太坊等EVM链上，机器人为了抢跑常设置极高的Gas价格；Solana中则表现为高优先费用（priority fee）。
    \item 交互模式规则：机器人通常只与少数核心协议（如DEX）交互，而人类用户可能使用更多类型的应用。
    \item 资金模式规则：机器人可能频繁从小型资金池中转入转出，形成特定资金流模式。
\end{enumerate}
\par 在《Detecting Financial Bots on the Ethereum Blockchain》中，作者设计了7条启发式规则（T1-T7）作为基线方法。实验表明，组合规则的召回率仅为43.8\%，远低于机器学习模型的77\%-80\%，说明简单规则难以捕捉复杂的机器人行为。
\par 启发式方法的局限性催生了基于机器学习的检测研究。

\subsection{基于机器学习的检测方法}

\par 机器学习方法通过从标注数据中自动学习行为特征与标签之间的关系，能够捕捉更复杂的模式，适应机器人策略的演化。根据标注数据的可用性，可分为监督学习、无监督学习和半监督学习。

\subsubsection {监督学习}
\par 监督学习需要一定规模的标注数据集（机器人/非机器人）。研究者从链上交易中提取多维特征，训练分类模型。常用模型包括随机森林（Random Forest）、梯度提升树（XGBoost、LightGBM）、支持向量机（SVM）和神经网络。
\par 特征工程是监督学习的核心。现有研究设计的特征可分为以下几类：
\begin{enumerate}
    \item 基础统计特征：交易总数、日均交易数、单日最大交易数、活跃天数等。
    \item 时间序列特征：交易间隔的均值与变异系数、每小时交易分布熵、夜间/周末交易比例、休眠模式特征（如GapBasedSleepiness）等。
    \item 交互多样性特征：交互的合约/程序种类数、交互代币种类数、程序调用熵、代币交易熵等。
    \item 经济特征：平均交易金额、gas费用统计、盈利/亏损模式等。
\end{enumerate}
\par 在《Detecting Financial Bots》中，作者提取了83维特征，涵盖上述类别，并采用随机森林（400棵树）进行分类，在二分类任务上取得83\%的准确率。特征重要性分析显示，与交易时间、频率和gas价格相关的特征（如Out-TX-Entropy、Out-TX-Per-Block、TX-GasPrice-Max）是区分机器人的关键。
\par 监督学习的优势在于准确率高，但依赖高质量标注数据。区块链领域缺乏公开的标注数据集，标注成本高，且可能存在标签噪声。

\subsubsection {无监督学习}
\par 无监督学习不依赖标签，通过聚类算法发现行为模式相近的地址簇，再通过人工解释簇的性质。常用算法包括K-Means、高斯混合模型（GMM）、DBSCAN等。
\par 在《Detecting Financial Bots》中，作者对比了K-Means和GMM在不同聚类数下的性能。最佳模型为GMM（30簇，均值填充，无降维），平均聚类纯度达到82.6\%。聚类结果揭示了不同类型的机器人簇，如CEX簇、Deposit Wallet簇、MEV簇等，验证了无监督学习的有效性。
\par 无监督学习的优势在于可以发现未知的机器人类型，但聚类结果需要人工解释，且对特征选择和参数敏感。此外，纯度和熵等内部评估指标不能完全保证聚类结果与真实类别一致。

\subsubsection {半监督学习与主动学习}
\par 在标签稀缺场景下，半监督学习利用少量标注数据和大量未标注数据提升性能。主动学习则通过迭代查询最不确定的样本进行人工标注，降低标注成本。目前这两种方法在区块链机器人检测中的应用尚少，但具有潜力。

\subsection{基于图神经网络的方法}

\par 区块链交易天然形成图结构（地址为节点，交易为边），图神经网络能够捕获地址间的资金流动模式和社区结构，为机器人检测提供新视角。
\par 研究者将地址交易图建模为有向加权图，节点特征包括交易统计量，边特征包括金额、时间等。通过图卷积网络、图注意力网络等模型学习节点表示，再用于分类或聚类。
\par 例如，Weber等人（2019）使用图卷积网络对比特币非法交易进行检测。在以太坊机器人检测方面，一些初步工作显示，结合交易图和时序特征的GNN模型能够提升检测性能。
\par 然而，GNN方法面临计算开销大、图构建依赖完整交易历史、可解释性弱等挑战，目前尚未在机器人检测领域广泛应用。

\subsection{多链环境下的机器人检测研究}

\par 随着多链生态的发展，研究者开始关注跨链机器人的行为模式以及检测方法在不同链上的迁移性。
\par 在BSC（BNB Chain）上，Chronis等人的系列研究（2023-2024）对狙击机器人进行了系统分析。他们通过构建流动性池数据集，识别狙击机器人的交易模式，并分析了反机器人机制的普及率。这些研究虽然聚焦于行为分析而非检测方法，但其特征提取思路（如交易速度、与代币部署时间的关系）可为检测研究提供参考。
\par 在Solana上，目前尚无系统的机器人检测研究。一些社区项目（如Solana-MEV-Bot-Blackbook）收集了疑似机器人地址列表，但缺乏规范的检测方法和验证。Solana的特殊架构——并行执行、历史证明（PoH）、无公开内存池——对检测方法提出了新的挑战，也带来了创新的机会。



\section{Solana生态与Meme币市场研究}

\subsection{Solana技术架构概述与MEV特性}

\par Solana是一条以高性能为设计目标的区块链，其核心架构与EVM链存在根本性差异。历史证明（PoH）机制、Sealevel并行执行引擎和无公开内存池的设计，一方面极大提升了吞吐量，另一方面也使MEV的提取从“监控内存池”转向了“与领导者合作”或“优化交易排序”。此外，Solana的本地费用市场（Local Fee Markets）使得针对特定账户的优先费用会随竞争程度波动，这为识别机器人提供了新的数据维度。


\subsection{Solana DeFi生态与核心协议}

\par Solana的高性能吸引了大量DeFi项目部署。Raydium作为最大的AMM DEX之一，同时集成了订单簿功能，其流动性池是机器人套利和狙击的主要目标。Orca采用集中流动性模型，交易执行效率更高，同样是机器人的高频交互对象。Jupiter聚合了多个DEX的流动性，为套利机器人提供了统一的交易接口。用户前期实验发现，大多数筛选出的机器人地址均与这三者深度交互，程序调用种类集中在5种左右——这种专注性恰恰体现了Solana机器人的专业化特征。
\par Jito Labs的出现更具启示性。这个类似Flashbots的MEV基础设施允许搜索者通过支付小费激励验证者将交易打包在有利位置，创造了一个公开的竞价市场。这意味着部分MEV活动从“隐蔽抢跑”转向了“公开竞价”，机器人行为也因此留下了更清晰的可追溯痕迹——优先费用的高低和波动性，有望成为区分普通用户与专业机器人的有效特征。

\subsection{Meme币在Solana的兴起与市场操纵研究}

\par Meme币是以玩笑或社区文化驱动的加密货币，缺乏实际应用价值，却凭借社交媒体传播实现短期内价格剧烈波动。Solana的低费用高速度，使其成为Meme币发行的热门平台。启动平台如pump.fun简化了代币创建和流动性建立流程，新代币可在几分钟内上线交易。这种快速启动模式，为狙击机器人创造了绝佳的狩猎场。
\par 由于Solana无公开内存池，狙击机器人无法像在以太坊上那样监控待处理交易。它们转而监控新代币的创建事件（如pump.fun的合约部署交易）或流动性池的初始化交易，一旦检测到新池建立，立即发送抢先购买交易，并设置高优先费用以确保被领导者优先打包。Meme币的高度投机性导致新代币上线初期交易量巨大、价格波动剧烈，这种环境天然适合机器人的快进快出策略。用户前期实验观察到，筛选出的活跃地址在新代币交互比例上存在明显分化，狙击型机器人往往具有更高的“新代币交互比例”，这一特征有望成为区分机器人类型的有效指标。
\par 论文《A Midsummer Meme's Dream》（arXiv'25）对Meme币市场的实证分析揭示了“人工增长 → 吸引投资者 → 利润提取”的完整操纵链条。这一发现表明，研究Solana机器人检测必须结合Meme币市场的操纵特征，将“与代币生命周期的关系”纳入特征设计。



\section{现有研究不足与本文切入点}

\par 通过对上述文献的系统梳理，可以发现当前研究存在以下空白与不足：
\begin{enumerate}
    \item 研究链路的空白：现有MEV与机器人行为研究高度集中于以太坊和BSC等EVM兼容链，而对Solana这一高性能非EVM链的关注严重不足。Solana独特的架构（无内存池、并行执行、本地费用市场）使得现有结论和方法难以直接迁移。
    \item 检测方法的局限：现有检测方法多依赖EVM链的公开内存池数据和Gas费用特征。针对Solana架构的机器人行为特征体系尚未建立，例如如何利用优先费用竞价痕迹、指令级程序调用信息以及新代币生命周期参与度来区分机器人，缺乏系统性的研究。
    \item 研究视角的单一：现有研究多将机器人视为同质化的“检测对象”，或仅聚焦于“狙击机器人”等单一类型。对于Solana生态中可能存在的不同类型机器人（如狙击型、套利型、做市型）及其行为分型，缺乏深入的分析与探讨。
\end{enumerate}
\par 基于上述分析，本文的切入点与创新设计如下：
\par 本文不简单移植现有EVM链的检测方法，而是立足Solana架构特性，构建一套适配Solana的机器人账户检测框架。重点从以下三个层面体现创新性：
\begin{enumerate}
    \item 凝练Solana特有特征及其行为含义：在特征工程中，除继承时间规律、交互多样性等经典维度外，新增优先费用的均值与变异系数、新代币交互比例、指令级程序调用熵等Solana特有指标。这些特征直接对应机器人在本地费用市场中的竞价行为、对Meme币生命周期的参与程度以及多指令策略的复杂性，使检测模型能够捕捉Solana生态特有的行为模式。
    \item 从“是否机器人”走向“机器人类型行为分型”：通过无监督聚类与监督学习相结合的策略，不仅识别地址是否为机器人，更进一步探索机器人的行为亚型（如狙击型、套利型、做市型）。通过分析不同类型机器人的特征差异，揭示其在Solana生态中的功能分化，为后续更精细化的监管与市场分析提供依据。
    \item 通过消融实验与案例研究增强方法说服力：在模型评估中，设计消融实验，量化Solana特有特征对检测性能的贡献，验证其必要性。同时，选取典型机器人地址进行深度案例分析，追踪其交易历史、资金流向及对特定Meme币价格的影响，将检测结果与链上实际行为相互印证，增强研究的实证支撑。
\end{enumerate}
\par 此外，本研究将探索引入LLM辅助标注与行为模式解读的可能性，利用LLM的语义理解能力，对聚类结果中的异常地址进行自动化解释，提升数据标注的效率与可解释性，为机器人检测方法提供新的技术路径。


\newpage
\begingroup
    \linespreadsingle{}
    \printbibliography[title={参考文献}]
\endgroup
